\documentclass[11pt]{article}
\usepackage[a4paper,margin=27mm]{geometry}
\usepackage{amsmath,amssymb,amsthm}
\usepackage[hidelinks]{hyperref}
\newtheorem{theorem}{Theorem}[section]
\newtheorem{lemma}[theorem]{Lemma}
\newtheorem{definition}[theorem]{Definition}
\newcommand{\Fix}{\operatorname{Fix}}
\newcommand{\id}{\mathrm{id}}
\newcommand{\K}{\mathrm K}
\newcommand{\llbracket}{[\![}
\newcommand{\rrbracket}{]\!]}
\title{Rigid projections, dependent sections, and represented universes}
\author{}
\date{11 November 2026}
\begin{document}
\maketitle

We develop dI-domain semantics, construct a represented reflexive domain
from a seed, and interpret the annotated term syntax in that domain. A Scott domain
is a pointed, bounded-complete algebraic dcpo. For the constructive
formulation, it has a small compact basis, a small-valued order, and
suprema of small directed and small bounded families. All joins below
are understood with this size restriction.

\section{Rigid projections}

Binary meets exist in a Scott domain: $x\wedge y$ is the bounded join of
the basic compact elements below both $x$ and $y$. A stable map is a
Scott-continuous map preserving meets of compatible pairs. The stable
order is
\[
 f\le_s g\quad\Longleftrightarrow\quad
 \forall x\le y,\quad f(x)=f(y)\wedge g(x).
\]
In particular this implies the pointwise inequality $f\le g$.

\begin{definition}
A rigid projection on $D$ is a continuous map $p:D\to D$ such that
$p\le_s\id$. Equivalently, $p(x)\le x$, and
\[
 x\le p(y)\quad\Longrightarrow\quad x\le p(x).
\tag{H}
\]
\end{definition}
The equivalence uses monotonicity. In the reverse direction, $p(x)$ is
below $p(y)\wedge x$ for $x\le y$, whereas this meet is fixed by (H) and
therefore below $p(x)$. In the forward direction, put $x=y$ to obtain
deflationarity and then apply the stable-order equation to $x\le p(y)\le y$.

\begin{theorem}
Every rigid projection is idempotent. Its fixed set is downward closed,
and $p(x)$ is the greatest fixed point below $x$. It preserves all binary
meets, including meets of incompatible elements. For rigid projections
$p,q$,
\[
 pq=qp=p\wedge q,
\]
where the last meet is pointwise. Their composition is again a rigid
projection and has fixed set $\Fix(p)\cap\Fix(q)$.
\end{theorem}
\begin{proof}
Apply (H) to $p(y)\le p(y)$ for idempotence. If $x\le y=p(y)$,
apply (H) again. If $y=p(y)\le x$, monotonicity gives $y\le p(x)$.
The element $p(x)\wedge p(y)$ is fixed, by downward closure, and is below
$x\wedge y$; the opposite inequality for meet preservation follows
by monotonicity. Finally, $p(q(x))$ is below both $p(x)$ and $q(x)$.
Their meet is fixed by $p$ and below $q(x)$, proving the reverse inequality.
Composition is continuous; its hereditary law follows from those of
$p$ and $q$.
\end{proof}

\section{Images and dI-domains}

\begin{theorem}
The fixed image of a rigid projection on a Scott domain is a Scott
domain. Its existing joins agree with ambient joins, and its compact
elements are exactly the ambient compact fixed points.
\end{theorem}
\begin{proof}
If $u=\bigvee_i x_i$ and each $x_i$ is fixed, then
$x_i=p(x_i)\le p(u)$, so $u\le p(u)\le u$. This also deals with
the empty join. For a fixed $x$, every basic compact below $x$ is fixed.
Consequently the ambient directed compact approximation of $x$ is an
approximation inside the image. Ambient compact fixed points remain
compact since directed joins are inherited. Conversely, a compact
element in the image equals one of the ambient basic compact elements
in its directed approximation. This proves both algebraicity and the
compactness characterization.
\end{proof}

A classical dI-domain is a Scott domain satisfying bounded
distributivity and the I-axiom: the lower set of each compact element is
finite. Here bounded distributivity means
\[
 x\wedge(y\vee z)=(x\wedge y)\vee(x\wedge z)
\]
whenever $y,z$ have a common upper bound.

\begin{theorem}
The fixed image of a rigid projection on a dI-domain is a dI-domain.
\end{theorem}
\begin{proof}
Meets and bounded joins in the image are the ambient ones, so
distributivity is inherited. A compact element of the image is compact
in the ambient domain. Its entire ambient lower set lies in the image,
so its finite lower-set presentation is inherited unchanged.
\end{proof}

The literal I-axiom requires care constructively. A finite discrete
presentation of every compact lower set, together with all small
directed suprema and a compact element $c\ne\bot$, implies excluded
middle for small propositions. Indeed, for a proposition $P$, take the
directed join $x_P$ of $\bot$ and a copy of $c$ indexed by $P$.
Then $x_P\le c$, so finite discreteness decides $x_P=c$. Compactness
and $c\ne\bot$ give $x_P=c$ if and only if $P$. Thus a constructive
event presentation with finite causal histories must not silently be
identified with this literal finite-discrete I-axiom. The image theorem
is conditional on the ambient dI structure; it does not construct a
nontrivial example under the constructive foundations alone.

\section{Constructing a projection from its fixed set}

\begin{theorem}\label{thm:subdomain}
Let $Z\subseteq D$ be a small-valued predicate that is downward closed
and closed under bounded joins, including the empty join. There is a
unique rigid projection with fixed set $Z$, namely
\[
 r_Z(x)=\bigvee\{k\mid k\text{ is basic compact},\ k\in Z,\ k\le x\}.
\]
Equivalently it is the greatest element of $Z$ below $x$.
\end{theorem}
\begin{proof}
The displayed family is small and bounded by $x$; closure places its
join in $Z$. If $y\in Z$ and $y\le x$, every basic compact below $y$
occurs in the family, proving $y\le r_Z(x)$. This gives monotonicity,
deflationarity and the hereditary law. To prove continuity, let $(x_i)$
be directed. A basic compact $k\le r_Z(\bigvee_i x_i)$ belongs to $Z$
and lies below some $x_i$, hence below $r_Z(x_i)$. Algebraicity gives
the required inequality. The greatest-element property proves uniqueness.
\end{proof}

For such predicates $Z,W$, inclusion gives $r_Z\le r_W$, and
\[
 r_{Z\cap W}=r_Zr_W.
\]
Thus constructing the fixed set first avoids any need to guess a
pointwise projector formula.

\section{The stable function domain}

\begin{theorem}\label{thm:stable}
For dI-domains $D,E$, the stable maps $D\to E$, ordered by the stable
order, form a Scott domain with a small compact basis. Directed joins
and bounded joins are pointwise. A stable map is compact in the stable
order if and only if it is compact in the ordinary continuous function
space with its pointwise order.
\end{theorem}
\begin{proof}
Stability and stable comparison can be tested on the small compact
bases. For example, compatible basic inputs have a common basic compact
upper bound. For general compatible inputs, continuity supplies basic
input witnesses for any compact output observation; directedness supplies
the common compact upper bound needed to apply the basic-input test.
The same method extends the stable-order equation from comparable basic
inputs to arbitrary $x\le y$. Thus both predicates are small-valued.

Below a common stable upper bound $h$, pointwise comparison implies
stable comparison. Indeed, if $f,g\le_s h$ and $f\le g$ pointwise, then
$f(y)\wedge g(x)\le f(y)\wedge h(x)=f(x)$ for $x\le y$; the other
inequality follows from monotonicity and $f\le g$.

Binary meets in any Scott domain preserve directed joins in each
argument, by compact approximation. Finite distributivity in $E$ extends
to all bounded joins: express a bounded join as the directed join of
its finite subjoins. These facts show that pointwise directed joins of
stable maps are stable, and that a pointwise join of maps bounded in
the stable order is stable and is their stable supremum. For example,
if $f_i,f_j\le_s h$ and $x,y\le z$, then
\[
 f_i(x)\wedge f_j(y)\le h(x\wedge y),\qquad
 f_i(x)\wedge f_j(y)\le f_i(x\wedge y).
\]
The second inequality follows from the first and the restriction equation
for $f_i\le_s h$. Distributing meets over the joins proves stability.

Every element below a compact element of a dI-domain is compact. To
justify this under the explicit finite-discrete formulation, fix a
finite presentation of $\mathord\downarrow c$. Equality in that lower
set is decidable. An element $x\le c$ either equals bottom, or gives
$c\ne\bot$ and hence local excluded middle by the argument above.
In the latter case, a directed family in a finite discrete poset has
a greatest member. Apply this to the directed basic approximation of
$x$, whose values lie in $\mathord\downarrow c$. Thus $x$ equals a
basic compact. The argument eliminates finite presentations and compact
witnesses only into propositions; it needs no choice of presentations.

For compact $a\in D$, $b\in E$ and stable $f$, put
\[
 k_{a,b,f}(x)=f(x\wedge a)\wedge b.
\]
Restriction to a principal lower ideal is a rigid projection, so this
map is stable and $k_{a,b,f}\le_s f$. It is compact in the ordinary
continuous function space. To see this, enumerate
$\mathord\downarrow a=\{u_1,\ldots,u_m\}$. A continuous map $g$ is
pointwise above $k_{a,b,f}$ exactly when
\[
 k_{a,b,f}(u_j)\le g(u_j)\quad(1\le j\le m).
\]
For sufficiency use $x\wedge a=u_j\le x$ and
$k_{a,b,f}(x)=k_{a,b,f}(u_j)$. Each output in these finitely many tests
is compact because it lies below $b$. Compactness of the outputs and
directedness give a common index satisfying all the tests.

These maps detect every compact observation of $f$: if compact
$b\le f(x)$, continuity supplies compact $a\le x$ with $b\le f(a)$,
and then $b\le k_{a,b,f}(x)$. Finite bounded joins of the approximations
are again stable and ordinarily compact. Ordinary compact maps have
finite compact-input/output graph presentations. Selecting the stable
ones gives a small basis, since stability is small-valued. This basis
is directed below every $f$. It determines both the pointwise order
and the stable restriction equation: for a compact observation below
$f(y)\wedge g(x)$ choose a basic stable approximation detecting it at
$y$, and apply that approximation's stable comparison with $g$.

Ordinary compactness of a stable map implies stable compactness: a
pointwise directed witness becomes a stable witness by the common-upper-bound
observation above. Conversely, stable compactness selects one element
of the directed basis approximation; hence the map equals an ordinarily
compact basis element.
\end{proof}

\begin{theorem}
The stable function space is a dI-domain. For compatible stable functions,
its meet is pointwise.
\end{theorem}
\begin{proof}
The pointwise meet of two stable functions is continuous and stable.
If $f,g\le_s h$, the restriction equations show that their pointwise
meet also lies below $h$ in the stable order. Pointwise comparison below
$h$ then proves the two lower-bound inequalities and the greatest-lower-bound
property.

For distributivity, write $u=y\vee z$ and $m=x\wedge u$. All of $m,y,z$
lie below $u$, so their compatible meets and bounded joins can be
computed pointwise. Distributivity in the codomain gives
$m=(m\wedge y)\vee(m\wedge z)=(x\wedge y)\vee(x\wedge z)$.
This argument permits $x$ to be incompatible with $u$.

For the I-axiom, first take a compact basis function with finite graph
$(a_j,b_j)_{j\in J}$, so
$f(x)=\bigvee_{a_j\le x}b_j$. Every $g\le_s f$ is reconstructed by
\[
 g=\bigvee_{j\in J} k_{a_j,c_j,f},\qquad
 c_j=g(a_j)\wedge b_j.
 \tag{R}
\]
Indeed, $c_j\le g(a_j)$ and the stable restriction equation imply
$k_{a_j,c_j,f}\le_s g$. Conversely, distribute $g(x)\wedge f(x)=g(x)$
over the active graph outputs. For $a_j\le x$, the restriction equation
and $b_j\le f(a_j)$ give
$g(x)\wedge b_j\le g(a_j)\wedge b_j=c_j\le k_{a_j,c_j,f}(x)$.
Thus the bounded join in (R) is exactly $g$.

Each $c_j$ lies in the finite discrete lower ideal of $b_j$. Enumerating
the finite product of these ideals and taking the corresponding joins
of restrictions gives a finite surjection onto $\mathord\downarrow f$.
Equation (R) also shows that equality of subfunctions is decided by the
finitely many clipped outputs. Duplicate removal therefore gives a
finite discrete presentation. The finite presentations of the output
ideals are used under propositional truncation; only finite choice,
proved by induction on the finite shape, is needed. Finally, any compact
stable function equals a compact basis function, so it inherits such a
presentation. This proof introduces no additional classical principle.
\end{proof}

\begin{theorem}
Binary products of dI-domains are dI-domains. Evaluation
\[
 [D\Rightarrow_s E]\times D\longrightarrow E,\qquad (f,x)\longmapsto f(x),
\]
is stable and continuous. For any Scott parameter domain $P$, currying
and uncurrying give mutually inverse operations between stable maps
$P\times D\to E$ and stable maps $P\to[D\Rightarrow_s E]$.
Their beta and eta laws are identities of stable maps.
\end{theorem}
\begin{proof}
Products have componentwise order, joins and meets, with products of
compact bases. Their distributivity is componentwise, and a compact
lower ideal is a product of two finite discrete compact lower ideals.

Continuity of evaluation follows from the pointwise directed-join law
and a common directed index for the function and argument observations.
For stability, let $f,g\le_s h$ and $x,y\le z$. Put
$t=f(x)\wedge g(y)$. Stability of $h$ gives $t\le h(x\wedge y)$.
The restriction equations for $f,g\le_s h$ then give
$t\le f(x\wedge y)$ and $t\le g(x\wedge y)$. The pointwise formula for
compatible function meets proves the required inequality; monotonicity
gives the other direction.

For a stable $H:P\times D\to E$, its sections $H(p,-)$ are stable.
If $p\le q$ and $x\le y$, stability applied to $(p,y),(q,x)\le(q,y)$
gives the restriction equation
$H(p,x)=H(p,y)\wedge H(q,x)$. Thus currying is monotone into the stable
order. Pointwise directed joins prove continuity, and stability applied
to $(p,x),(q,x)$ proves preservation of compatible meets in the parameter.
Uncurrying is stable by pairing and evaluation. Extensionality proves
beta and eta.
\end{proof}

\section{Dependent sections over a fixed domain}

For dI-domains $A,E$, take the stable function domain
$F=[A\Rightarrow_s E]$ from Theorem~\ref{thm:stable}. Its order is the stable
order; its bounded joins are pointwise. Let $x\mapsto b_x$ be a
monotone family of rigid projections on $E$, and set
\[
 S_b=\{f\in F\mid \forall x\in A,\ b_x(f(x))=f(x)\}.
\]
The following construction only uses that $F$ is a Scott domain of
continuous maps, its order implies pointwise order, and its bounded
joins are pointwise.

\begin{lemma}
Membership in $S_b$ can be tested at basic compact inputs. In particular
it is small-valued when the input basis is small.
\end{lemma}
\begin{proof}
Suppose $b_a(f(a))=f(a)$ for basic compact $a$. For $a\le x$,
monotonicity of the family gives $f(a)\le b_x(f(a))\le b_x(f(x))$.
Continuity of $f$ and the directed compact approximation of $x$ yield
$f(x)\le b_x(f(x))$. Deflationarity gives equality.
\end{proof}

\begin{theorem}
$S_b$ is a downward-closed, bounded-join-closed subdomain of $F$.
Its rigid projection is
\[
 R_b(f)=\bigvee\{k\mid k\text{ is basic compact in }F,
                         \ k\le_s f,\ k\in S_b\}.
\]
For two families $b,c$, write $(b\cap c)_x=b_xc_x$. Then
\begin{align*}
 b\le c&\ \Longrightarrow\ R_b\le R_c,\\
 R_{b\cap c}&=R_bR_c,\\
 R_{b\cap c}(f\wedge g)&=R_b(f)\wedge R_c(g).
\end{align*}
The meets of functions in the last line are in the order of $F$.
\end{theorem}
\begin{proof}
Stable order implies pointwise order, and each fibre is downward closed,
so $S_b$ is downward closed. Bounded joins are pointwise and each fibre
is closed under bounded joins. Apply Theorem~\ref{thm:subdomain}.
The first two laws follow from inclusion and intersection of the fixed
sets. For the last law, $R_b(f)\wedge R_c(g)$ is a section of both
families, by downward closure, and is below $f\wedge g$. Conversely,
any common section below $f\wedge g$ is below both $R_b(f)$ and $R_c(g)$.
\end{proof}

This differs essentially from pointwise clipping. On the two-element
chain, take $b_0=\bot$ and $b_1=\id$. Clipping the constant-one function
gives the identity, but $\id\not\le_s(\lambda x.1)$: the stable-order
equation at $0\le1$ fails. The actual rigid section projection sends
the constant-one function to bottom.

\section{Finite tests and parameter continuity}

The preceding theorems give the section domain for each family.
The finite compact presentations now prove continuity in the family.

The finite-test property is the following: for each compact
function $k$ there are finitely many compact input/output pairs
$(a_j,c_j)$ such that, for every monotone family $b$,
\[
 k\in S_b\quad\Longleftrightarrow\quad
 \forall j,\ b_{a_j}(c_j)=c_j.                 \tag{FT}
\]
The finite graph presentation
$k(x)=\bigvee\{c_j\mid a_j\le x\}$ supplies this property. For necessity
use $c_j\le k(a_j)$ and downward closure. For sufficiency, each active
$c_j$ belongs to the fibre at $x$, by monotonicity, and that fibre is
closed under the displayed bounded join. The presentations need not be
canonical minimal traces. They are available by Theorem~\ref{thm:stable}.

\begin{lemma}
For a directed family $(b_i)$ with pointwise supremum $b$,
\[
 R_b(f)=\bigvee_i R_{b_i}(f).
\]
\end{lemma}
\begin{proof}
The right side is below the left by monotonicity. If compact
$k\le R_b(f)$, then $k\in S_b$ and $k\le f$. Apply (FT). For each of its
finitely many tests,
$c_j=b_{a_j}(c_j)=\bigvee_i(b_i)_{a_j}(c_j)$.
Compactness gives an index passing that test. Directedness gives a
single index passing all tests, hence $k\in S_{b_i}$ and
$k\le R_{b_i}(f)$. Algebraicity completes the proof.
\end{proof}

More precisely, let $P$ be any Scott parameter domain, and suppose
$p\mapsto b^p$ is monotone and each map
$p\mapsto b^p_x(z)$ is continuous. Then $p\mapsto R_{b^p}(f)$ is
continuous. If $p\mapsto f_p$ is also continuous, so is
$p\mapsto R_{b^p}(f_p)$: apply continuity separately in the two arguments
and then use directedness to dominate both indices by a common index.
Both assertions are checked in Unit.

\begin{theorem}
Let $\mathcal R(E)$ denote the rigid projections on $E$. This is a
dI-domain, represented by the principal ideal below the identity in
$[E\Rightarrow_s E]$. For stable families in
$\mathcal F=[D\Rightarrow_s\mathcal R(E)]$, both maps
\[
 \mathcal F\times[D\Rightarrow_s E]\longrightarrow[D\Rightarrow_s E],
 \quad (b,f)\longmapsto R_b(f),
\]
\[
 \mathcal F\longrightarrow\mathcal R([D\Rightarrow_s E]),
 \quad b\longmapsto R_b
\]
are stable.
\end{theorem}
\begin{proof}
The principal ideal is a rigid image, so inherits the dI structure.
Every element below identity satisfies the rigid projection equation,
and conversely every rigid projection defines such an element.
All rigid projections have the common stable upper bound identity;
therefore their meets are pointwise and equal composition.
Compatible meets in $\mathcal F$ are consequently fibre intersections.

Fibre evaluations on this actual family domain are continuous, by
stable evaluation and the inherited pointwise directed joins in the
rigid image. The preceding parameter-continuity theorem therefore
applies to $(b,f)\mapsto R_b(f)$. Its joint intersection equation proves
preservation of compatible meets. Currying gives a stable map into the
stable endomorphism domain of $[D\Rightarrow_s E]$; each value is a rigid
projection, so it factors through the principal ideal below identity.
The formal computation theorem identifies this factor with the original
greatest-section construction.
\end{proof}

The earlier section construction and its intersection laws apply to
arbitrary monotone families. The stable family space above supplies the
parameter domain for the following varying-input construction.

\section{Varying rigid input domains}

For a rigid projection $p$ on $D$, define input restriction on the stable
function domain $F=[D\Rightarrow_s E]$ by $I_p(f)=f\circ p$.

\begin{theorem}
$I_p$ is a rigid projection on $F$, and $p\mapsto I_p$ is stable as a map
$\mathcal R(D)\to\mathcal R(F)$. Stable maps $\Fix(p)\to E$ correspond
by extension and restriction to ambient stable maps $f:D\to E$ satisfying
$f\circ p=f$.
\end{theorem}
\begin{proof}
Internal composition is stable, by evaluation, pairing and currying.
Since $p\le_s\id$, composition gives $f\circ p\le_s f$.
For heredity, suppose $g\le_s f\circ p$. Apply the restriction equation
at $p(x)\le x$. Idempotence of $p$ gives
\[
 g(p(x))=g(x)\wedge f(p(p(x)))=g(x),
\]
since $g(x)\le f(p(x))$. Thus $g\le_s g\circ p$, as required.
Currying internal composition and restricting its values to the
principal ideal below identity proves stability in $p$.

The inclusion $\Fix(p)\to D$ and projection $D\to\Fix(p)$ are stable;
their composite on the image is identity and their ambient composite
is $p$. Precomposition with these maps gives extension and restriction,
with the stated beta and eta laws.
\end{proof}

\begin{theorem}
For a stable family $B:D\to\mathcal R(E)$, put
\[
 Q_{p,B}=R_B I_p=I_p R_B.
\]
This is a rigid projection on $F$ and is stable in $(p,B)$, also jointly
with its function argument. $Q_{p,B}(f)$ is the greatest $g\le_s f$ such that
\[
 g\circ p=g,\qquad B_x(g(x))=g(x)\quad(x\in\Fix(p)).
\]
It depends only on the restriction of $B$ to $\Fix(p)$.
\end{theorem}
\begin{proof}
Rigid projections commute and their composition is their meet.
Composition of rigid projections is itself stable: restrict internal
composition to the two principal ideals below identity, and then
restrict its output to the same ideal. Composing this operation with
$p\mapsto I_p$ and $B\mapsto R_B$ proves the stability assertions.

For a supported $g$, the fibre condition at fixed inputs implies the
condition at every $x$: $g(x)=g(p(x))$ belongs to the fibre at $p(x)$,
and the family is monotone with $p(x)\le x$. Thus the two displayed
conditions say precisely that $g$ is fixed by both $I_p$ and $R_B$.
The greatest-fixed-point property of their composition gives the result.
It also proves independence of the family away from fixed inputs.
Finally, any stable family defined directly on $\Fix(p)$ extends along
$p$, so this ambient representation covers those families too.
\end{proof}

\section{Stable recursion and projector stages}

\begin{theorem}
For a dI-domain $D$, the least-fixed-point operation
$Y:[D\Rightarrow_s D]\to D$ is stable. Its finite stages are an increasing
chain in the stable function order. Stable parameter families of
endomorphisms consequently have stable least solutions.
\end{theorem}
\begin{proof}
Let $X=[D\Rightarrow_s D]$ and $H=[X\Rightarrow_s D]$. The functional
\[
 T:H\to H,\qquad T(g)(f)=f(g(f)),
\]
is stable by evaluation, pairing and currying. Its Kleene iterates
starting at the bottom of $H$ are increasing in the stable order.
Their directed join is an element of $H$, hence a stable map $X\to D$.
At each $f$, induction identifies the iterates with $f^n(\bot)$, and
pointwise directed joins identify their limit with the ordinary Kleene
least fixed point. Parameter stability follows by composition.
\end{proof}

In particular, a stable recursive operator on $\mathcal R(D)$ has rigid
stages $p_n$ and a rigid least solution $p_\infty$, with
$p_\infty(x)=\bigvee_n p_n(x)$. Since $p_m\le p_n$ for $m\le n$,
\[
 p_m p_n=p_m,\qquad p_n p_\infty=p_n.
\]
Stable feedback into the varying-family parameters, followed by
$(p,B)\mapsto Q_{p,B}$, is one such recursive projector operator.
These are abstract Kleene stages. No finiteness of the stages or
identification with the earlier Scott-model rank stages is claimed.

\section{A conditional internal universe}

A \emph{rigid embedding} $e:A\to D$ comes with a stable retraction
$r:D\to A$ satisfying $re=\id_A$ and $er\le_s\id_D$.
Then $er$ is a rigid projection and $e$ is strict. Transport of rigid
projections is stable:
\[
 \mathcal R(A)\longrightarrow\mathcal R(D),\qquad q\longmapsto eqr.
\]
Its action on embedded arguments is $eqr(e(x))=e(q(x))$.

The following is a theorem \emph{from explicit representation data}.
The next section constructs these data from an arbitrary seed dI-domain;
nontriviality of that seed remains an explicit requirement for a
nontrivial model.

\begin{theorem}
Suppose a dI-domain $D$ has rigid embeddings
\[
 e_U:\mathcal R(D)\longrightarrow D,\qquad
 e_F:[D\Rightarrow_s D]\longrightarrow D,
\]
with stable retractions $r_U,r_F$. There are a stable typing operation
$\Delta:D\times D\to D$, a universe element $U$ satisfying $\Delta(U,U)=U$,
and a stable dependent-function code constructor
$\Pi:D\times[D\Rightarrow_s D]\to D$.
\end{theorem}
\begin{proof}
Set $\Delta(a,x)=r_U(a)(x)$ and $u=e_Ur_U$, regarded as a rigid
projection on $D$. Let $U=e_U(u)$. Then $r_U(U)=u$, so
\[
 \Delta(U,x)=u(x),\qquad
 \Delta(U,e_U(q))=e_U(q),\qquad \Delta(U,U)=U.
\]
Conversely, $\Delta(U,a)=a$ implies $a=e_U(r_U(a))$, so universe members
are exactly the represented rigid projections. Each $\Delta(a,-)$ is
rigid; stability of $\Delta$ follows from rigid evaluation.

For $a\in D$ and stable $b:D\to D$, let $p=r_U(a)$ and
$B(x)=r_U(b(x))$. Define
\[
 \Pi(a,b)=e_U\bigl(e_F Q_{p,B}r_F\bigr).
\]
Every operation in this expression has been proved stable. The result
is a code, hence lies in $U$, and
\[
 \Delta(\Pi(a,b),e_F(f))=e_F(Q_{p,B}(f)).
\]
Stable term application is $\operatorname{app}(v,x)=r_F(v)(x)$, with
$\operatorname{app}(e_F(f),x)=f(x)$.
The greatest-supported-section theorem gives dependent-function
introduction, including abstraction via $f\circ p$ for functions
specified on fixed inputs. For elimination, if
$x\le\Delta(\Pi(a,b),x)$, applying $r_F$ shows
$r_F(x)\le_s Q_{p,B}(r_F(x))$. Thus $r_F(x)$ is a section on fixed inputs.
For $y\le\Delta(a,y)$, downward closure of the fibre gives
\[
 \operatorname{app}(x,y)\le
 \Delta(b(y),\operatorname{app}(x,y)).
\]
The inequalities express fixedness, since the typing operations are
deflationary. These introduction and elimination laws are checked.
\end{proof}

\begin{theorem}
The code embedding $e_U$ alone gives a stable recursive-type operation
on stable code-valued functionals $G:D\to\mathcal R(D)$. Its result
$\mu(G)$ lies in $U$ and satisfies
\[
 \mu(G)=e_U(G(\mu(G))).
\]
It is least among all $x\in D$ satisfying $e_U(G(x))\le x$.
Its stages start at bottom, lie in $U$, unfold by this equation, increase,
and have supremum $\mu(G)$.
\end{theorem}
\begin{proof}
Define $\widehat G(c)=G(e_U(c))$ on $\mathcal R(D)$ and set
$\mu(G)=e_U(Y(\widehat G))$. Internal composition and the stable
least-fixed-point operator prove stability in $G$. The fixed-point
equation gives unfolding, and the result lies in $U$ since it is a code.

If $e_U(G(x))\le x$, apply $r_U$ to get $G(x)\le r_U(x)$.
Since $e_U(r_U(x))\le x$, monotonicity gives
$\widehat G(r_U(x))\le r_U(x)$. Leastness on the code domain yields
$Y(\widehat G)\le r_U(x)$, hence $\mu(G)\le e_U(r_U(x))\le x$.
The stages are $e_U(\widehat G^n(\bot))$; strictness and continuity of
$e_U$ give their bottom start and limit, and every stage is a code.
\end{proof}

\section{A represented universe from a seed}

The code embedding does not need an independent construction.
The domain $\mathcal R(D)$ is the rigid principal image below identity
in $[D\Rightarrow_s D]$. Composing this image embedding with any rigid
embedding $[D\Rightarrow_s D]\to D$ supplies $e_U$.

\begin{theorem}
For every dI-domain $A$ there is a dI-domain $D$ with a rigid embedding
$A\to D$ and a stable isomorphism
\[
 D\cong[D\Rightarrow_s D].
\]
Consequently $D$ has the represented universe, stable dependent-function
codes and stable recursive-type operation of the previous section.
The seed embedding is faithful, so nontriviality of $A$ implies
nontriviality of $D$.
\end{theorem}
\begin{proof}
First, rigid embeddings lift through stable endomorphism spaces. If
$e:A\to B$ and $r:B\to A$ form a rigid embedding/retraction pair, the
lifted pair is
\[
 f\longmapsto efr,\qquad g\longmapsto rge.
\]
These maps are stable by internal composition. Their first composite is
identity. The second is $g\mapsto pgp$, where $p=er$, and is below
identity in stable order: both precomposition and postcomposition by
$p\le_s\id$ decrease a stable map. The hereditary property is checked,
so this is again a rigid embedding. There is also a starting embedding
$A\to[A\Rightarrow_s A]$, sending $a$ to the constant function $a$,
with retraction evaluation at bottom. Its composite is the rigid input
restriction $f\mapsto f\circ\bot$.

We need closure under countable rigid limits. For a chain
$(A_n,e_n,r_n)$, let
\[
 L=\{(x_n)\mid r_n(x_{n+1})=x_n\},
\]
ordered coordinatewise. Rigid retractions preserve bounded joins, so
both directed and bounded joins in $L$ are coordinatewise. Write
$\iota_n:A_n\to L$ for the canonical extension by embeddings to the
right and projections to the left, and $\pi_n:L\to A_n$ for evaluation.
They satisfy
\[
 \pi_n\iota_n=\id,\qquad
 \iota_n(a)\le x\ \Longleftrightarrow\ a\le\pi_n(x).
\]
Both maps are stable, and $p_n=\iota_n\pi_n$ is rigid. Moreover
$p_n\le_s p_{n+1}$ and $\sup_n p_n=\id_L$.

The elements $\iota_n(a)$, for stage basis elements $a$, form a small
compact basis. Compactness follows from the displayed adjunction;
directedness uses a common later stage and compactness there, and
density is coordinatewise. Distributivity transfers coordinatewise.
If $a$ is compact at a stage, the lower ideal of $\iota_n(a)$ is entirely
inside the image of $\iota_n$, by rigidity. A finite discrete
presentation of $\mathord\downarrow a$ therefore transfers to one of
$\mathord\downarrow\iota_n(a)$. Every compact of $L$ equals a basis
element, so $L$ satisfies the finite-ideal axiom as well.

The stable endomorphism construction commutes with this limit. For the
chain of endomorphism spaces with the lifted embeddings, the inverse
stable maps are
\[
 f\longmapsto(\pi_n f\iota_n)_n,
 \qquad
 (f_n)_n\longmapsto\sup_n\iota_n f_n\pi_n.
\]
The restrictions form a stable cone. The assembly maps increase in
stable order: the $n$th assembly equals the next one precomposed with
the $n$th rigid projection of the function-sequence limit. Hence their
supremum is stable, also as a function of the sequence.

For the first inverse law, assembly after restriction is
$\sup_n p_nfp_n=f$. To verify the nontrivial inequality, continuity
reduces to inputs $p_nx$ and density reduces outputs to
$p_mf(p_nx)$. Each is below the diagonal approximation at
$\max(m,n)$. For the other inverse law, compression to level $n$ of an
assembly at level $m\le n$ is below $f_n$, while compression of any
later assembly is exactly $f_n$. Coherence and absorption of nested
stage projections prove these claims. Continuity of compression then
gives the inverse equation.

Now set $A_0=A$ and $A_{n+1}=[A_n\Rightarrow_s A_n]$, starting with
the constant-function embedding and lifting each successive embedding.
Let $D$ be its rigid limit. The chain of endomorphism spaces is precisely
the tail of this tower. Deleting the first coordinate is a stable
isomorphism of limits, with inverse adjoining $r_0(x_1)$. Thus
\[
 [D\Rightarrow_s D]\cong\lim_n[A_n\Rightarrow_s A_n]
   =\lim_n A_{n+1}\cong D.
\]
The canonical embedding of $A_0$ gives the seed embedding. Its retraction
proves faithfulness. Composing the rigid-code image embedding with the
function-space isomorphism supplies the universe representation.
\end{proof}

This is a construction relative to any seed dI-domain, with no assumed
function-space or code embedding for the seed. It does not supply an
unconditional nontrivial seed under constructive foundations: the
finite-discrete compact-ideal obstruction discussed earlier still
applies. The existing Scott filter model has not been shown to provide
such a seed. The following sections supply a conditional concrete seed and connect
the resulting reflexive universe to terms and tower-stage observations.

\section{A conditional nontrivial Boolean seed}

Write $2=\{0<1\}$. The classical input in this section is an explicit
parameter
\[
 \mathsf{EM}_{\ell}:\quad P+\neg P
 \quad\text{for each proposition }P:\mathcal U_{\ell}.
\]
It is not a parameter of the ambient formal development.

\begin{theorem}
Under $\mathsf{EM}_{\ell}$, the two-element chain is a dI-domain with
both elements compact. Its associated represented reflexive domain is
nontrivial. Conversely, a Scott-domain structure on this Boolean order
with bottom $0$ and compact top $1$ implies $\mathsf{EM}_{\ell}$.
\end{theorem}
\begin{proof}
For a small family $(x_i)_{i\in I}$, decide the proposition
$\|\sum_{i:I}(x_i=1)\|$. Its join is $1$ when this proposition holds,
and $0$ otherwise. This gives all small joins, including the empty join.
If $1\le\bigvee_i x_i$, the definition supplies a mere index with
$x_i=1$, proving compactness of $1$. Compactness of $0$ uses the
inhabitedness of a directed family. The whole Boolean carrier is a
compact basis. Distributivity is a Boolean case distinction, and the
lower ideals have the explicit presentations $\{0\}$ and $\{0,1\}$.
Neither of these last two verifications needs excluded middle.

The seed embedding in the reflexive construction is faithful and
strict. Thus its two images are distinct, and the upper image is
nonbottom. For the converse, apply the compact-ideal excluded-middle
argument to $1$ and its two-element lower ideal.
\end{proof}

This supplies a concrete classical instance while exposing exactly the
logical input used. It does not claim a nontrivial literal dI-domain
under the constructive global assumptions alone.

\section{Stable interpretation of annotated syntax}

Let $D$ have the representation data constructed above. Put
$E_0=1$ and $E_{n+1}=E_n\times D$. A term with $n$ free variables is
interpreted as a stable map $E_n\to D$. For an environment $\rho$, set
$a=\llbracket A\rrbracket\rho$ and
$f(x)=\llbracket M\rrbracket(\rho,x)$. The clauses are
\begin{align*}
 \llbracket U\rrbracket\rho&=U,\\
 \llbracket\Pi(A,B)\rrbracket\rho
   &=\Pi\bigl(a,x\mapsto\llbracket B\rrbracket(\rho,x)\bigr),\\
 \llbracket\lambda(A,M)\rrbracket\rho
   &=e_F\bigl(f\circ r_U(a)\bigr),\\
 \llbracket MN\rrbracket\rho
   &=\operatorname{app}(\llbracket M\rrbracket\rho,
                          \llbracket N\rrbracket\rho).
\end{align*}
Variables are product projections. The restriction in the lambda clause
is needed for the supported-section interpretation of dependent functions.

\begin{theorem}
These clauses give a total stable interpretation of the existing
\texttt{Syntax.Tm}. It respects renaming and simultaneous substitution.
For an argument $x=\llbracket N\rrbracket\rho$,
\[
 \llbracket(\lambda(A,M))N\rrbracket\rho
   =\llbracket M\rrbracket(\rho,\Delta(a,x)).
\]
Consequently, if $\Delta(a,x)=x$, this equals
$\llbracket M[N]\rrbracket\rho$.
\end{theorem}
\begin{proof}
All constructors are stable by composition, pairing and currying;
input restriction is jointly stable in the projection and function.
Renaming is proved by induction on terms, extending the environment
comparison with an identical last argument under each binder.
Extensionality of stable maps identifies the curried bodies. The same
induction proves substitution; weakening, obtained from renaming,
handles substitution under a binder. The represented application beta
law gives the displayed equation, and single substitution gives the
last assertion.
\end{proof}

\begin{theorem}
Write $t:a$ at $\rho$ for
$\llbracket t\rrbracket\rho\le
\Delta(\llbracket a\rrbracket\rho,\llbracket t\rrbracket\rho)$.
If $M:B$ at every $(\rho,x)$ with $\Delta(a,x)=x$, then
$\lambda(A,M):\Pi(A,B)$ at $\rho$. If
$f:\Pi(A,B)$ and $N:A$ at $\rho$, then $fN:B[N]$ at $\rho$.
Also $U:U$, every Pi expression lies in $U$, and beta holds at typed
arguments.
\end{theorem}
\begin{proof}
Introduction is the greatest-supported-section abstraction theorem;
elimination is its application theorem, followed by semantic single
substitution in the result type. Deflationarity turns semantic typing
into fixedness, giving the beta premise.
\end{proof}

These are semantic rules for the new interpretation. Full induction over
\texttt{Syntax.Der}, including Pi congruence and typed function
extensionality, is not yet formalised. In particular, the old Scott-model
soundness and adequacy results are not being reused as results about this
different interpretation.

\section{Tower rank and finite observation}

For any rigid chain limit, let $p_n=\iota_n\pi_n$ and say that $x$ has
rank at most $n$ when $p_nx=x$.

\begin{theorem}
Rank is downward closed and persists at later stages. Every compact
observation $k\le x$ is captured at some finite stage:
\[
 \exists n,\quad k\le p_nx.
\]
Every compact element has a finite-rank witness. For any stable map
$f:P\to D$, the maps $p_n\circ f$ form an increasing chain in stable
order with supremum $f$. In particular this applies to every term
interpretation and every environment.
\end{theorem}
\begin{proof}
Downward closure follows from rigidity, and persistence from
$p_n\le_s p_m$ for $n\le m$. Since $x=\bigvee_n p_nx$, compactness
supplies the observation witness. Heredity then gives $p_nk=k$.
Postcomposition preserves stable comparison. Pointwise directed joins
and $\bigvee_n p_n=\id$ identify the supremum of the composed maps.
The existential witnesses are propositionally truncated; no least rank
or simultaneous choice of ranks is asserted.
\end{proof}

The exponential tower gives a more specific link to application. If
$v=(v_n)$ is an element of its reflexive limit and $r_F(v)$ is its decoded
stable function, then
\[
 \pi_n\,r_F(v)\,\iota_n=v_{n+1}.
\]
This is the $n$th coordinate of the exponential-limit inverse equation,
since dropping the first coordinate of $v$ gives its function sequence.
Consequently the following equation is checked:
\[
 p_n\bigl(\operatorname{app}(v,p_nx)\bigr)
   =\iota_n\bigl(v_{n+1}(\pi_nx)\bigr).
\]
It connects a finite input/output observation of application to the next
function stage, without appealing to the old Scott-filter syntax ranks.

For the Boolean seed, rank zero already contains two distinct points.
These stages are therefore not the bottom-starting Kleene iterations.
We have not proved finiteness of the whole Boolean tower stage carriers,
a bound on observations from syntactic term rank, or an identification
with the earlier Scott-model stages or Proposition~2.27.

\section{What is formalised}

The following files are included in the combined Unit check:
\begin{itemize}
\item \texttt{ScottMeet.seg}: binary meets from the small compact basis.
\item \texttt{RigidProjector.seg}: the equivalence of the hereditary law
and the stable-order equation, idempotence, downward closure, meet
preservation, composition and commutation.
\item \texttt{RigidImage.seg}: the Scott-domain image and both directions
of the compactness characterization.
\item \texttt{DIDomain.seg}, \texttt{DIImage.seg}: an explicit abstract
dI-domain interface and inheritance of distributivity and finite
discrete compact lower ideals by rigid images. Finite presentations
are propositionally truncated, so no simultaneous choice is required.
\item \texttt{RigidSubdomain.seg}: construction from a small
downward-closed bounded-join-closed predicate, the converse, round trip,
monotonicity and intersection.
\item \texttt{RigidSections.seg}: compact-input tests, the section
projection, its greatest-section property, monotonicity, intersection,
and the joint intersection equation, over an explicitly supplied
function-domain interface.
\item \texttt{StableMaps.seg}, \texttt{StableDirected.seg},\\
\texttt{DistributiveJoins.seg}, \texttt{StableComplete.seg}: small tests
for stability and stable order, and pointwise directed and bounded joins.
\item \texttt{FiniteDiscrete.seg}, \texttt{CompactIdeals.seg}: finite
discrete positions, locally derived excluded middle, and compactness
of elements below compact elements.
\item \texttt{StableComposition.seg}, \texttt{PrincipalProjection.seg},\\
\texttt{StablePatches.seg}: stable composition, principal restrictions,
and the ordinarily compact stable approximations $k_{a,b,f}$.
\item \texttt{StableExponential.seg}: the Scott domain of stable maps,
its small compact basis, and equivalence of stable and ordinary
compactness for stable maps.
\item \texttt{StableMeet.seg}, \texttt{StableDistributive.seg},\\
\texttt{StableFiniteIdeals.seg}: pointwise compatible meets, distributivity,
and finite discrete compact lower ideals; the stable exponential as a
\texttt{DIDomain.DI}.
\item \texttt{FiniteEnumeration.seg}, \texttt{FiniteRange.seg}: finite
products, mere finite choice, decidable tests and duplicate removal,
with small indices even for large enumerated types.
\item \texttt{ScottProduct.seg}, \texttt{DIProduct.seg},
\texttt{StableProduct.seg}: products and stable projections and pairing.
\item \texttt{StableEvaluation.seg}, \texttt{StableCurrying.seg}:
joint stable evaluation, currying and uncurrying, and beta/eta identities.
\item \texttt{StableSections.seg}: the genuine stable function-domain
instance of the section construction and the finite tests (FT) on its
compact basis.
\item \texttt{SectionContinuity.seg}: continuity in any Scott parameter
with continuous fibre evaluations, including a simultaneously varying
continuous function argument.
\item \texttt{RigidSpace.seg}, \texttt{StableFamilies.seg}: the dI-domain
of rigid projections, the stable family domain, joint stability of section
application, and stability of the family-to-projection operation.
\item \texttt{StableInternalComposition.seg}, \texttt{RigidOperations.seg}:
internal composition and rigid composition/evaluation as stable maps.
\item \texttt{StableInputRestriction.seg}, \texttt{StableImageMaps.seg},\\
\texttt{StableVaryingFamilies.seg}: input restriction, image extension and
restriction, and the stable greatest-supported-section construction.
\item \texttt{StableFixpoint.seg}, \texttt{StableProjectorFixpoint.seg}:
stable least-fixed-point formation and rigid projector stages and limits.
\item \texttt{RigidEmbedding.seg}, \texttt{RigidUniverse.seg}:
explicit rigid representation data, stable transport, and the conditional
universe, application and dependent-function typing theorems.
\item \texttt{RigidUniverseRecursion.seg}: stable recursive-type formation,
leastness, and stages, conditional on the code representation.
\item \texttt{RigidEmbeddingComposition.seg}, \texttt{StableConjugation.seg},\\
\texttt{RigidEmbeddingExponential.seg}, \texttt{RigidEmbeddingLaws.seg}:
embedding constructors, stable conjugation, exponential lifting,
compactness and bounded-join preservation.
\item \texttt{RigidLimitOrder.seg}, \texttt{RigidLimitEmbedding.seg},\\
\texttt{RigidLimitDomain.seg}, \texttt{RigidLimitRigid.seg},
\texttt{RigidLimitDI.seg}: countable rigid limits, their compact basis,
rigid stage embeddings and both dI axioms.
\item \texttt{OrderedChains.seg}, \texttt{StableExtensional.seg},\\
\texttt{RigidLimitLaws.seg}, \texttt{RigidLimitMaps.seg}: stage identities,
identity approximation and stable assembly of cones.
\item \texttt{RigidLimitExponential.seg},
\texttt{RigidLimitExponentialLaws.seg},\\
\texttt{RigidLimitTail.seg}: the stable exponential limit isomorphism
and the tail isomorphism, with both inverse laws.
\item \texttt{RigidExponentialTower.seg}, \texttt{RigidReflexiveDomain.seg}:\\
the relative reflexive-domain construction, faithful seed embedding,
represented universe and recursive types.
\item \texttt{BooleanDI.seg}, \texttt{BooleanReflexive.seg}: the conditional
Boolean seed, its necessary classical input, and two distinct points
in its represented reflexive domain.
\item \texttt{ScottTerminal.seg}, \texttt{StableSyntax.seg},\\
\texttt{StableSubstitution.seg}, \texttt{StableSyntaxTyping.seg}:\\
stable term interpretation, renaming, substitution, and introduction,
elimination and beta laws for dependent functions.
\item \texttt{RigidRankStages.seg}, \texttt{ReflexiveSyntax.seg}: tower ranks,
compact-observation capture, stable term approximation, and the
one-step application equation.
\end{itemize}

A nontrivial seed without a classical parameter, full judgemental
soundness for the stable interpretation, comparison with syntactic and
Scott-model ranks, and a constructive event-structure presentation
remain open. The conditional Boolean instance takes small excluded
middle explicitly; the general dI theorems continue to take only their
stated domain and representation data.
The formal development still uses only the existing univalence and
replacement interfaces as global assumptions. The combined check
reports 7,745 declarations and zero failures, with all 40 frozen
dependency hashes verified.

\end{document}
